\documentclass[11pt,a4paper]{article}
\usepackage{fontspec}
\usepackage{kotex}
\usepackage{geometry}
\usepackage{graphicx}
\usepackage{xcolor}
\usepackage{tcolorbox}
\usepackage{booktabs}
\usepackage{array}
\usepackage{longtable}
\usepackage{colortbl}
\usepackage{fancyhdr}
\usepackage{titlesec}
\usepackage{hyperref}
\usepackage{emoji}
\usepackage{amsmath}
\usepackage{amssymb}
\usepackage{enumitem}

\geometry{left=25mm, right=25mm, top=30mm, bottom=30mm}

\setmainfont{Noto Sans CJK KR}
\setsansfont{Noto Sans CJK KR}
\setmonofont{Noto Sans CJK KR}
\setemojifont{Noto Color Emoji}

\definecolor{primary}{HTML}{1976D2}
\definecolor{darkbrown}{HTML}{3E2723}
\definecolor{mediumbrown}{HTML}{5D4037}
\definecolor{lightbrown}{HTML}{8D6E63}
\definecolor{cream}{HTML}{FFF3E0}
\definecolor{lightgray}{HTML}{F5F5F5}
\definecolor{tableheader}{HTML}{E8E8E8}

\titleformat{\section}
  {\Large\bfseries\color{primary}}
  {\thesection.}{0.5em}{}
\titleformat{\subsection}
  {\large\bfseries\color{darkbrown}}
  {\thesubsection}{0.5em}{}
\titleformat{\subsubsection}
  {\normalsize\bfseries\color{mediumbrown}}
  {\thesubsubsection}{0.5em}{}

\renewcommand{\contentsname}{목차}
\renewcommand{\figurename}{그림}
\renewcommand{\tablename}{표}

\tcbset{
  tipbox/.style={colback=cream,colframe=primary,boxrule=1.5pt,arc=3mm,
    left=5mm,right=5mm,top=3mm,bottom=3mm,fonttitle=\bfseries\color{primary}},
  formulabox/.style={colback=lightgray,colframe=lightgray,boxrule=0pt,arc=2mm,
    left=5mm,right=5mm,top=3mm,bottom=3mm},
  summarybox/.style={colback=white,colframe=primary,boxrule=2pt,arc=4mm,
    left=5mm,right=5mm,top=5mm,bottom=5mm}
}

\hypersetup{colorlinks=true,linkcolor=primary,urlcolor=primary,bookmarks=true,unicode=true}


\pagestyle{fancy}
\fancyhf{}
\fancyhead[R]{\small\color{lightbrown}수업 집중도 분석 프로그램 이론해설서 v1.0.0.0}
\fancyfoot[C]{\thepage}
\renewcommand{\headrulewidth}{0.4pt}

\begin{document}

\begin{titlepage}
\centering
\vspace*{3cm}
{\Huge\bfseries\color{primary}\emoji{bar-chart} 수업 집중도 분석 프로그램}\\[0.5cm]
{\Huge\bfseries\color{primary} 이론해설서}
\\[1.5cm]
{\Large 시뮬레이터 버전: v2.4.0.0}
\\[2cm]
\begin{tcolorbox}[summarybox]
\centering
{\large 이 이론해설서는 수업 집중도 분석 프로그램이 다루는 핵심 개념과 원리를 설명합니다.}
\end{tcolorbox}
\vfill
{\large 사물인터넷과 빅데이터}\\[0.3cm]
{\large 청주교육대학교 컴퓨터교육과}
\vspace{1cm}
\end{titlepage}

\section*{1. 이 문서의 목적}

이 이론해설서는 \textbf{수업집중도 분석 프로그램}이 다루는 핵심 개념과 원리를 설명합니다. 시뮬레이터를 사용하기 전에 배경 지식을 쌓거나, 사용 후 깊이 있는 이해를 위해 참조합니다.


\begin{tcolorbox}[summarybox]
\textbf{핵심 주제}: 상관과 인과, 회귀 분석

\end{tcolorbox}


\section*{2. 핵심 개념}

\begin{center}
\renewcommand{\arraystretch}{1.4}
\begin{tabular}{|p{2.5cm}|p{12.5cm}|}
\hline
\rowcolor{tableheader}
\textbf{개념} & \textbf{설명} \\
\hline
\textbf{상관관계} & 두 변수가 함께 변하는 경향. 양의 상관(함께 증가), 음의 상관(하나 증가 시 다른 하나 감소) \\
\hline
\textbf{인과관계} & 한 변수의 변화가 다른 변수의 변화를 직접 일으키는 관계. 상관 ≠ 인과 \\
\hline
\textbf{회귀 분석} & 데이터에서 변수 간 관계를 수식(y = ax + b)으로 표현하는 통계 기법 \\
\hline
\textbf{R² (결정계수)} & 모델이 데이터를 얼마나 잘 설명하는지의 비율 (0\textasciitilde{}1, 1에 가까울수록 좋음) \\
\hline
\textbf{최소자승법} & 예측값과 실제값의 차이(잔차) 제곱합을 최소화하는 계수를 찾는 방법 \\
\hline
\textbf{규칙 기반 모델} & 사람이 직접 기준을 설정하는 if-else 방식의 예측 모델 \\
\hline
\end{tabular}
\end{center}


\section*{3. 원리 상세}

\subsection*{상관 ≠ 인과}

\textbf{핵심 원칙}: "데이터는 원인을 알려주지 않고, 우리가 원인을 추정한다."


아이스크림 판매량과 익사 사고가 양의 상관을 보입니다. 아이스크림이 익사를 유발할까요? 아닙니다. 공통 원인(여름 더위)이 둘 다 증가시킬 뿐입니다. 이처럼 상관관계가 있다고 인과관계가 있는 것은 아닙니다.


\subsection*{규칙 기반 vs 회귀 모델}

\begin{center}
\renewcommand{\arraystretch}{1.4}
\begin{tabular}{|p{2.2cm}|p{6.5cm}|p{6.2cm}|}
\hline
\rowcolor{tableheader}
\textbf{구분} & \textbf{학생 모델 (규칙)} & \textbf{교수 모델 (회귀)} \\
\hline
기준 설정 & 사람이 직접 & 데이터에서 자동 계산 \\
\hline
방식 & if 온도 > 28 then 집중도↓ & y = 0.3x₁ - 0.2x\textsubscript{2} + 2.1 \\
\hline
장점 & 이해하기 쉬움 & 더 정확 (데이터 기반) \\
\hline
단점 & 주관적, 편향 가능 & 해석 어려움 \\
\hline
핵심 & "내가 정한 기준" & "데이터가 찾은 패턴" \\
\hline
\end{tabular}
\end{center}

두 모델 모두 "AI가 판단하는 것"이 아니라 "사람이 설계한 규칙을 실행"하는 것입니다.



\section*{4. 실생활 응용 사례}

\textbf{사례 1}: 학교: 교실 환경(온도·소음·CO\textsubscript{2})과 학습 효과의 관계 연구


\textbf{사례 2}: 병원: 환자 활력 징후(체온·혈압·심박)와 회복 속도의 관계 분석


\textbf{사례 3}: 공장: 설비 온도·진동과 불량률의 관계 모니터링


\textbf{사례 4}: 마케팅: 광고 비용·노출 빈도와 매출의 관계 회귀 분석



\section*{5. 더 알아보기}

$\bullet$ \textbf{다중 회귀 분석}: 변수가 3개 이상일 때의 회귀 모델

$\bullet$ \textbf{교차 검증}: 모델의 일반화 성능을 평가하는 기법

$\bullet$ \textbf{과적합}: 학습 데이터에만 잘 맞고 새 데이터에 약한 문제

$\bullet$ \textbf{설명 가능한 AI (XAI)}: AI의 판단 근거를 사람이 이해할 수 있게 만드는 분야



\section*{6. 핵심 용어 사전}

\begin{center}
\renewcommand{\arraystretch}{1.4}
\begin{tabular}{|p{2.3cm}|p{2.1cm}|p{10.6cm}|}
\hline
\rowcolor{tableheader}
\textbf{용어} & \textbf{학술 용어} & \textbf{쉬운 설명} \\
\hline
상관관계 & — & 두 변수가 함께 변하는 경향 \\
\hline
인과관계 & — & 한 변수의 변화가 다른 변수의 변화를 직접 일으키는 관계 \\
\hline
회귀 분석 & — & 데이터에서 변수 간 관계를 수식(y = ax + b)으로 표현하는 통계 \\
\hline
R² (결정계수) & R² (결정계수) & 모델이 데이터를 얼마나 잘 설명하는지의 비율 (0\textasciitilde{}1, 1에 가까울수록 \\
\hline
최소자승법 & — & 예측값과 실제값의 차이(잔차) 제곱합을 최소화하는 계수를 찾는 방법 \\
\hline
규칙 기반 모델 & — & 사람이 직접 기준을 설정하는 if-else 방식의 예측 모델 \\
\hline
\end{tabular}
\end{center}


\section*{7. 참고자료}

\begin{center}
\renewcommand{\arraystretch}{1.4}
\begin{tabular}{|p{1.5cm}|p{6.9cm}|p{6.6cm}|}
\hline
\rowcolor{tableheader}
\textbf{\#} & \textbf{자료명} & \textbf{비고} \\
\hline
1 & 수업집중도분석프로그램\_퀵가이드 & 소프트웨어 사용 중 빠른 참조 \\
\hline
2 & 수업집중도분석프로그램\_사용설명서 & 전체 기능 상세 \\
\hline
3 & 수업집중도분석프로그램\_활용가이드 & 수업 적용 \\
\hline
4 & 수업집중도분석프로그램\_개발참조 & 기술 사양 \\
\hline
\end{tabular}
\end{center}


\vfill
\begin{center}
\small\color{lightbrown}
\textit{수업 집중도 분석 프로그램 이론해설서}
\end{center}
\end{document}